\section{Implementation}
\label{sec:implementation}
We implement \sys as a Python framework consisting of around  8K lines of code.
We implement a custom C\&C server and use open-source malware  capabilities (from the Caldera project~\cite{caldera}) to infect and send shell commands to hosts.\footnote{We picked Caldera given our familiarity.
Other C\&C servers such as  Cobalt Strike~\cite{cobalt} or Merlin~\cite{merlin} could also be used.}
We implement \sys with custom Python modules.
For the environment state service, we create parsers that interpret command outputs and update the knowledge base.

For each of the five high-level tasks in \secRef{sec:high_level_api}, we create non-LLM and LLM-based agents that translate the tasks into low-level primitives (e.g., Python scripts, shell commands).
We implement the non-LLM agents for the lateral movement and privilege escalation tasks by integrating into \sys's internal library (or optionally Metasploit's library) of known vulnerabilities and their corresponding exploits.
For instance, if an LLM specifies to lateral move into a host with a CVE, \sys\ will identify the CVE and execute the low-level exploit.


We use LangChain~\cite{LangChain2025} to iteratively prompt LLMs.
We first create a prompt with the onboarding process outlined in \secRef{sec:high_level_api}.
During the execution phase, we extract the Python function between the \texttt{<task></task>} or \texttt{<query></query>} tags.
Then \sys executes the function to get a list of tasks for the orchestrator to execute.
\sys will execute attacks until an LLM specifies a \texttt{<finished>} tag or reaches a time limit.

\para{\benchmark} To systematically evaluate \sys, we design and implement \benchmark, a multi-host red teaming benchmark with 40 environments. We use Python and Ansible code to set up the environments atop OpenStack.
The red team  goal in  10 environments is to exfiltrate key data files and the goal in the remaining  30 environments is to gain root access to key hosts.
We design \benchmark to be diverse along key dimensions: (1) network size and topology, (2) type of vulnerabilities, and (3) red teaming complexity. 

In terms of network size,  \benchmark 
currently includes many  small enterprise environments ranging from 22 to 50 hosts.
For 30 of the environments, we algorithmically generate different topology structures that are similar to real-world environments~\cite{ciscoEnterpriseNetwork, ibmTreeNetwork}.
In addition, we manually design the topology of 10 environments based on topologies from prior work such as ``Star'', ``Chain'', and ``Dumbbell''~\cite{mirage, ringNetwork2013technique, ferguson2021_deception_psychology, ciscoEnterpriseNetwork, dumbbellNetwork} and topologies from public reports of real-world attacks~\cite{equifax_report, colonial_pipeline_techtarget}.
The manually designed topologies are named based on the environment they were adapted from (e.g., Equifax environment). The   algorithmically generated topologies are named based on the topology structure: ``N4-H41-G7'' has four (sub)networks, 41 hosts and 7 critical assets (i.e., goals).

In terms of vulnerabilities, \benchmark includes diverse vulnerabilities such as common misconfigurations (e.g., plain text credentials), remote code execution vulnerabilities (e.g., Apache Struts CVE-2017-5638), and privilege escalation vulnerabilities (e.g., \texttt{sudo} CVE-2021-3156).
Several of these vulnerabilities have been used in real-world attacks~\cite{equifax_ftc_settlment, colonial_pipeline_techtarget}.

In terms of red teaming complexity,   we vary the number of critical assets as well as the attack graph complexity. Across the environments, red team success  spans  a spectrum ranging  from  2 to 48 assets and from 5 to 104 tasks.
More details on the specifics of each environment are in \appRef{sec:appendix_environments}.

